% Response to the AI review
% Same text as rebuttal/ai_review.md (1491 characters as posted).
\section*{Response to the AI Review}

Thank you for a precise review and for valuing our cost accounting and failure-regime reporting; we agree with most points.

\textbf{Attribution:} Agreed. Our only matched-budget control, p.~5's interpolated host point, favors the host by $\sim$1.5 mAP, and no random-skip control was run. We will state PSM's own share (60.7 $\rightarrow$ 46.6 GFLOPs/frame, 8 of the 73 points vs.\ dense), drop the ``three quarters'' reading and claim composability, not superiority.

\textbf{Recurrence:} Agreed: both protocols test reuse within short clips after a bootstrap, not after long gaps. We will scope the persistence claim; gap, cache age and drift are untested.

\textbf{Complexity:} Agreed. Skipped tokens keep the host's K/V (Alg.~1, l.~14), so savings are linear in $\gamma k$; we will rederive Eqs.~3--4 and the break-even for this path with one FLOP convention. Table 1 uses op counters, not these equations; the counters omit matching ($\sim$0.25 GFLOPs/frame), which we will add.

\textbf{Novelty:} We will cite MC-ViT, VidToMe and LITE and present PSM as a training-free, content-addressed reuse policy, not a new primitive.

\textbf{Latency:} We claim no speed-up: PSM is near host latency and behind the other gated methods, and at 1024 spatial pooling is faster and more accurate (pp.~5--6).

\textbf{Recognition:} We will report the dense gap beside every K400 result and narrow those claims.

\textbf{Minor:} All accepted: the 49\% figure, Fig.~1 caption, query wording, frozen memory, subsample overlap, and no reading within noise.

Thank you again!
